\introduction

La sostenibilidad se define como la capacidad de satisfacer las necesidades del presente sin comprometer la capacidad de las futuras generaciones para satisfacer sus propias demandas, abarcando la interrelación entre el desarrollo económico, la inclusión social y la protección del medio ambiente \parencite{WCED1987}. Este concepto implica un enfoque holístico que busca un equilibrio entre estos tres pilares para garantizar un futuro viable para todos. Las dimensiones de la sostenibilidad se dividen comúnmente en tres categorías principales: la dimensión económica, centrada en el crecimiento y la eficiencia de recursos a largo plazo \parencite{Stiglitz2018}; la dimensión social, que promueve la equidad, la inclusión y el bienestar humano \parencite{NacionesUnidas2015}; y la dimensión ambiental, enfocada en la conservación de la biodiversidad y la gestión responsable de los recursos naturales \parencite{EuropeanUnion2024}.

La educación para el desarrollo sostenible es fundamental en la promoción de un futuro equitativo. Este tipo de educación no solo busca transmitir conocimientos, sino que también fomenta habilidades y actitudes que permiten a las personas tomar decisiones informadas y responsables. La Educación para el Desarrollo Sostenible (EDS) es un enfoque integral que empodera a los ciudadanos para participar activamente en la construcción de sociedades justas y sostenibles \parencite{UNESCO2017, UNESCO2020EDS}. Además, la EDS contribuye directamente a la consecución de los Objetivos de Desarrollo Sostenible (ODS), especialmente a la meta 4.7 del ODS 4, garantizando que las personas adquieran los conocimientos teóricos y prácticos necesarios para promover estilos de vida sostenibles \parencite{NacionesUnidas2015, UNESCO2020EDS}.

La educación es un derecho humano fundamental y un motor clave para el desarrollo sostenible. Sin embargo, persisten brechas críticas: en 2021, 244 millones de niños y jóvenes estaban sin escolarizar, cifra que refleja desigualdades arraigadas en geografía y nivel socioeconómico \parencite{GlobalEduReport2022}. Estas brechas se intensifican en crisis humanitarias y se ven exacerbadas por fenómenos globales que limitan el acceso a la educación a distancia \parencite{UNICEF2020}. En América Latina, la falta de dispositivos y conectividad afecta significativamente la continuidad educativa \parencite{BancoMundial2025}. Por ello, abordar estas brechas resulta esencial para lograr una educación inclusiva, equitativa y de calidad para todos.

En Cuba, el sistema educativo constituye históricamente un pilar del desarrollo social. El 99{,}7\% de su población adulta está alfabetizada \parencite{UNESCO2023}, y la educación es gratuita en todos los niveles tal como lo establece la Constitución de la República \parencite{CubaConstitucion2019}. En el panorama nacional, el desarrollo guiado por ingeniería formal de software interactivo para los ODS discurre en su fase embrionaria, aunque cimentado en proyecciones tecnológicas provenientes del ámbito universitario y apuestas sólidas como el plan de informatización escolar Aula Digital. El archipiélago orienta gradualmente el perfeccionamiento continuo del capital lúdico digital introducido a las aulas de educación básica escolarizada \parencite{InformeVoluntarioCuba2021}. No obstante, persisten desafíos como la escasez de materiales didácticos digitales, la obsolescencia tecnológica y las limitaciones de conectividad que afectan la calidad educativa. Las fluctuaciones de suministro en el ecosistema físico acentúan un rezago en la disponibilidad masiva de periféricos de última línea, mermando el campo efectivo de las iniciativas puramente conectadas y exigiendo arquitecturas optimizadas \parencite{BancoMundial2023, MesaLago2022}.

Los videojuegos se posicionan como herramientas pedagógicas innovadoras para abordar desafíos de sostenibilidad. Desde una perspectiva histórica, los juegos de mesa constituyen herramientas pedagógicas por excelencia durante siglos, ofreciendo entornos estructurados donde los participantes desarrollan habilidades cognitivas, sociales y estratégicas de manera lúdica. La digitalización de estos juegos representa una evolución natural que mantiene la esencia pedagógica del formato original mientras amplifica sus posibilidades mediante la interactividad digital, la retroalimentación inmediata y la eliminación de barreras logísticas \parencite{MartinezGimenes2023}. Esta transición del formato físico al digital no implica una simple traslación de mecánicas, sino una reconfiguración que aprovecha las capacidades computacionales para enriquecer la experiencia de aprendizaje.

En el ámbito internacional, las Naciones Unidas desarrollan el juego de mesa \textit{Go Goals!} como herramienta pedagógica oficial para introducir los diecisiete ODS a niños de entre ocho y diez años \parencite{GoGoals2020}. Este juego utiliza una estructura inspirada en el clásico Serpientes y Escaleras, adaptando sus mecánicas tradicionales para incorporar contenidos educativos específicos sobre sostenibilidad mediante un sistema de preguntas y respuestas que refuerzan el aprendizaje. Su simplicidad y familiaridad facilitan la comprensión de las reglas, permitiendo que el foco cognitivo se centre en el contenido educativo.

El conocimiento de los ODS conlleva múltiples beneficios: alineación con la Agenda Global, mejora en la toma de decisiones y fomento de la responsabilidad social y ambiental \parencite{R3crea2022}. Por el contrario, su desconocimiento genera desalineación con tendencias internacionales, uso inadecuado de recursos y menor preparación para enfrentar desafíos globales como el cambio climático \parencite{EcologiaVerde2024}.

A pesar de su reconocimiento internacional, el formato físico del juego \textit{Go Goals!} presenta limitaciones estructurales inherentes que restringen su acceso, escalabilidad, rejugabilidad y capacidad de evaluación del aprendizaje. A esto se suma que la adaptación digital específica de juegos de mesa reconocidos internacionalmente como \textit{Go Goals!} representa un vacío identificado en la oferta actual de recursos educativos tecnológicos en Cuba. No se documenta la existencia de un videojuego desarrollado en el país que se centre exclusivamente en la enseñanza de los ODS mediante esta mecánica lúdica particular. Esta ausencia evidencia una brecha digital significativa considerando la creciente asimilación de los videojuegos en el público joven y su demostrada capacidad para facilitar el Aprendizaje Basado en Juegos. Aprovechando el potencial pedagógico de los videojuegos educativos y la estructura probada del juego de mesa original, se identifica la necesidad de desarrollar una versión digital que supere las barreras del formato físico y mitigue esta brecha en el contexto educativo.

A partir de la situación problemática antes expuesta, se plantea como problema de la investigación: ¿Cómo contribuir a la enseñanza de los Objetivos de Desarrollo Sostenible? Como objeto de estudio se tienen los videojuegos educativos como herramientas para la promoción de los ODS, enmarcados en el campo de acción: el proceso de desarrollo de un videojuego educativo basado en la adaptación digital del juego de mesa \textit{Go Goals!}. El objetivo general es desarrollar un videojuego educativo basado en la adaptación digital del juego de mesa \textit{Go Goals!} que contribuya a la promoción de los Objetivos de Desarrollo Sostenible.

\subsection*{Tareas de la investigación}

\begin{itemize}
    \item Sistematización de los referentes teóricos y metodológicos que sustentan el uso de videojuegos educativos para la promoción de los Objetivos de Desarrollo Sostenible.
    \item Análisis de soluciones homólogas y diagnóstico del estado actual del juego analógico \textit{Go Goals!} para determinar los elementos susceptibles de informatización.
    \item Selección de las tecnologías, herramientas y metodologías idóneas para la implementación del sistema.
    \item Elaboración del diseño arquitectónico, definición de mecánicas y especificación de requisitos funcionales y no funcionales del videojuego.
    \item Implementación del videojuego educativo a partir de la adaptación digital del referente analógico seleccionado.
    \item Validación del sistema mediante la ejecución de pruebas de calidad de software que avalan su despliegue y pertinencia.
\end{itemize}

\subsection*{Métodos teóricos de investigación}

Los métodos teóricos de investigación resultan esenciales para explorar y comprender fenómenos complejos en distintas disciplinas académicas. Permiten analizar, interpretar y sintetizar teorías y conceptos, fomentando el avance del conocimiento y la generación de nuevas perspectivas. A través del análisis conceptual, la revisión bibliográfica y la síntesis teórica, los investigadores profundizan en las bases teóricas de sus estudios, estableciendo conexiones significativas entre ideas y enfoques diversos. Estos métodos no solo fundamentan la investigación, sino que también contribuyen a la construcción de conocimiento sólido en el ámbito académico.

Analítico-sintético: se utiliza para extraer las ideas fundamentales que respaldan teóricamente la investigación y la propuesta presentada. Este enfoque se justifica debido a la necesidad de establecer una base sólida y coherente para el estudio, asegurando que las ideas esenciales sean identificadas y comprensiblemente integradas en el marco teórico. La revisión exhaustiva de la literatura existente permite situar el estudio dentro del contexto académico pertinente, proporcionando un fundamento robusto para el desarrollo y la argumentación de la investigación propuesta.

Modelación: constituye una herramienta esencial en la investigación, adoptando formas tanto tangibles como abstractas, con el propósito de replicar y comprender a fondo el objeto de estudio. Funciona como una representación concreta o conceptual de un fenómeno o proceso, permitiendo una mayor claridad en la comprensión de sus características, desarrollo y funcionamiento. A través de esquemas, diagramas o representaciones, se capturan las interrelaciones y propiedades clave de la realidad, facilitando así la exploración y el análisis detallado de la materia investigada.

\subsection*{Métodos empíricos}

Los métodos empíricos en la investigación constituyen un enfoque fundamental que se basa en la observación y la experiencia directa para obtener datos y validar hipótesis. Estos métodos implican la recopilación sistemática de información a través de la experimentación, la observación o la recolección de datos cuantitativos y cualitativos en el entorno real. Al emplear métodos empíricos, los investigadores pueden poner a prueba teorías, validar resultados y generar conocimiento sustentado en evidencia concreta y verificable.

Observación: se emplea para analizar soluciones homólogas existentes en el ámbito de los videojuegos educativos sobre sostenibilidad, identificando características, mecánicas de juego y enfoques pedagógicos que informan el diseño de la solución propuesta.

Entrevista: facilita una comprensión más profunda de los requerimientos del cliente al interactuar con la aplicación. Esto ayuda a recopilar información directamente del cliente y, durante la fase de pruebas, a evaluar su satisfacción con los resultados obtenidos.

\subsection*{Estructura capitular}

\textbf{Capítulo 1}: Fundamentación teórica, en el cual se examinan en detalle los conceptos fundamentales que comprenden el marco teórico de la investigación. Se ahonda, además, en el ámbito de los videojuegos educativos y sus características. También se estudian soluciones homólogas para establecer comparaciones. Asimismo, se exploran las herramientas, tecnologías que se emplean para el desarrollo del sistema y las metodologías de desarrollo de software.

\textbf{Capítulo 2}: Diseño e implementación, donde se detalla la propuesta de solución a través de la metodología de desarrollo de software seleccionada. Se definen los requisitos funcionales y no funcionales del sistema y se detallan las interfaces de usuario y tarjetas CRC para documentar la parte visual del proyecto y estructurar la arquitectura de la solución.

\textbf{Capítulo 3}: Resultados y validación, en el que se documentan los resultados de las pruebas realizadas en el despliegue de la solución.
